%!TEX root = ../main.tex
% SÍMBOLOS matemáticos o de notación usados en el documento.
%
% Definir:   \newglossaryentry{sym:clave}{type=symbols, name={\ensuremath{...}},
%                description={...}, sort={clave}}
% Usar:      \gls{sym:clave}

\newglossaryentry{sym:e}{
    type=symbols,
    name={\ensuremath{E}},
    description={Exposición de un riesgo, producto de su probabilidad por su impacto},
    sort={e}
}

\newglossaryentry{sym:f1}{
    type=symbols,
    name={\ensuremath{F_1}},
    description={Media armónica entre precisión y exhaustividad},
    sort={f1}
}

\newglossaryentry{sym:g}{
    type=symbols,
    name={\ensuremath{G = (V, E)}},
    description={Grafo dirigido con vértices \ensuremath{V} y aristas \ensuremath{E}},
    sort={g}
}

\newglossaryentry{sym:skf}{
    type=symbols,
    name={\ensuremath{S, K, F}},
    description={Conjuntos de nodos source, sink y sanitizer del análisis de contaminación},
    sort={s}
}

\newglossaryentry{sym:tp}{
    type=symbols,
    name={\ensuremath{TP, FP, TN, FN}},
    description={Verdaderos positivos, falsos positivos, verdaderos negativos y falsos negativos},
    sort={tp}
}
